\section{Experimental Setup}
\label{sec:experimental_setup}

\subsection{Testbed and serving layer}
\label{ssec:exp_testbed}

All experiments use the \boptest{} \texttt{bestest\_air} testcase, a
single-zone air-side HVAC reference building. \boptest{} is served through
the \texttt{boptest\_rte} container, which exposes the standard \boptest{}
HTTP API to both surrogate training and live closed-loop validation.
The \texttt{bestest\_air} envelope is a single thermal zone
($48\,\mathrm{m}^{2}$ floor area, lightweight construction following the
ANSI/ASHRAE 140 BESTEST 600 baseline), heated by an idealised air-side
unit that accepts a supply-temperature command in
$[20,60]\,^{\circ}\mathrm{C}$. Weather is drawn from the \boptest{}
default TMY file for Brussels, Belgium, with a $15$-minute control
interval and a comfort band of $[21,24]\,^{\circ}\mathrm{C}$ throughout
the heating season. The Block~3 hydronic testcases share this
single-zone envelope class but differ in actuator interface; the
testcase manifest is summarised in Table~\ref{tab:architecture} of
Section~\ref{ssec:res1_architecture} and detailed in
Section~\ref{sec:results_transfer}.

\subsection{Datasets and sample generation}
\label{ssec:exp_datasets}

Three transition corpora support the experiments. The
\texttt{v3} hourly corpus is the original logging frequency of the
direct-supply-temperature trajectories and is used only to bootstrap
the data-driven surrogate; the \texttt{v3.5} prepared $15$-minute
corpus is the canonical training set after the Stage~A preprocessing
of Section~\ref{ssec:method_v35}; and the \texttt{v3.5} collected
$15$-minute corpus is the raw form of the same trajectories before
preprocessing, retained for ablation. Supplementary Table~S1 lists the
provenance, parameter ranges, and physical rationale of each corpus,
and Supplementary Table~S2 documents the numerical justification of
the chosen sample sizes (including the learning-curve check that
fixes the canonical corpus at $\sim\!2{\times}10^{5}$ transitions).
The canonical artefact for the calibrated twin is the prepared
$15$-minute corpus; the hourly corpus is used only for the
\texttt{v3} pre-training run referenced in
Section~\ref{ssec:method_v3}.

\subsection{Preprocessing, scaling and feature engineering}
\label{ssec:exp_preprocessing}

Preprocessing follows the Stage~A pipeline documented in
Supplementary Table~S3 (latency compensation, bias removal, power
affine normalisation, rolling denoise, causal $\Delta T$
recomputation). Per-channel scaling parameters and the reverse
transform used at inference time are listed in Supplementary
Table~S7. Pairwise input independence is verified in Supplementary
Table~S5 through Pearson correlation and mutual-information checks;
no input pair exceeds the documented thresholds.

\subsection{Evaluation Protocol}
\label{ssec:exp_eval_protocol}

The evaluation protocol is fixed across all controller families and
\textbf{never varies} between Section~\ref{sec:results_fidelity} and
Section~\ref{sec:results_control}.

\paragraph{Scenarios.}
Two live \boptest{} closed-loop windows are used as the primary transfer
benchmarks:
\begin{itemize}
    \item \texttt{peak\_heat\_window}: simulated start
    $t_{\mathrm{start}} = 14$~days after $1$~January, duration
    $14$~days, capturing the coldest contiguous window of the Brussels
    TMY file.
    \item \texttt{typical\_heat\_window}: simulated start
    $t_{\mathrm{start}} = 108$~days after $1$~January, duration
    $14$~days, capturing a milder average heating period.
\end{itemize}
Yearly validation uses the full \boptest{} year with KPI windows defined by
the \boptest{} specification.

\paragraph{Simulation step and comfort band.}
The simulation step is fixed at $\Delta t = \SI{900}{\second}$ (\SI{15}{\minute}).
The comfort band is $[T_{\mathrm{low}}, T_{\mathrm{high}}] = [\SI{21}{\celsius}, \SI{24}{\celsius}]$.

\paragraph{Metrics.}
\begin{itemize}
    \item $\msmetric$: the \boptest{} composite KPI combining thermal
    discomfort and energy use; lower is better.
    \item \texttt{violation\,\%}: fraction of simulation steps spent outside
    the comfort band.
    \item \texttt{energy}: total final energy consumption in kWh.
    \item \texttt{RMSE\_T}: root-mean-square error of zone temperature
    against a reference trajectory; reported separately for one-step
    replicative validity and for multi-horizon predictive rollout.
    \item \texttt{first\_divergence\_step}: the first simulation step at
    which the surrogate-trained policy's trajectory diverges from the live
    \boptest{} trajectory by more than a fixed threshold (transfer-realism
    diagnostic).
\end{itemize}

\paragraph{Statistical reporting.}
\begin{itemize}
    \item Predictive validity numbers (Section~\ref{sec:results_fidelity})
    report bootstrap $95\%$ confidence intervals over rollout windows,
    inherited directly from
    \texttt{outputs/surrogate\_v35\_rollout\_prepared\_15min\_power\_head\_only/calibrated\_v35/horizon\_metrics.csv}.
    \item Canonical controller results (Section~\ref{sec:results_control})
    are reported as \mbox{$\mathrm{mean} \pm \mathrm{std}$} over
    $N=3$ random seeds (\texttt{42}, \texttt{43}, \texttt{44}) for the
    thermostatic hybrid, the pure \texttt{v3} baseline, and the MORL
    canonical configuration. HDRL is reported single-seed; this is
    acknowledged as a limitation in Section~\ref{sec:discussion}.
\end{itemize}

\subsection{Reproducibility statement}
\label{ssec:exp_reproducibility}

All code, configuration, trained models, and the eleven supplementary
tables are released through the project repository and the associated
archival deposit; the URL and DOI are listed on the manuscript
metadata page. The pre-registered audit chain comprises four
git-anchored commits:
(a)~MORL canonical-selection pre-registration
(\texttt{93df9b3});
(b)~MORL post-$N{=}5$ canonical falsification
(\texttt{62dc859});
(c)~Block~3 transferability pre-registration
(commit recorded in the manifest field
\texttt{audit.pre\_registration\_commit\_sha} of
\texttt{configs/block3\_testcase\_manifest.yaml}); and
(d)~Block~3 closure
(\texttt{audit.block3\_close\_commit\_sha}).
The \boptest{} version is \texttt{0.7.0}; the testcase identifiers
are \texttt{bestest\_air}, \texttt{bestest\_hydronic\_heat\_pump},
\texttt{bestest\_hydronic}, and
\texttt{singlezone\_commercial\_hydronic}. The software stack is
Python~$3.11$, PyTorch~$2.1$, NumPy~$1.26$, and Stable-Baselines3
$2.2$ \citep{Raffin2021SB3}. All experiments run on a single CPU
thread (Intel Core~i7-class, $32$~GB RAM, no GPU required); random
seeds are fixed in \texttt{configs/seeds.yaml}. Numerical values in
Tables~\ref{tab:architecture}--\ref{tab:control} are reproduced
automatically from the corresponding CSVs in \texttt{reports/} by
\texttt{paper/build\_paper\_tables.py}, so every cell in the main
tables is traceable to a deterministic source file.

\subsection{Runtime characteristics}
\label{ssec:exp_runtime}

A throughput comparison between the live \boptest{} RTE HTTP loop and
the three surrogate backends is reported in Table~\ref{tab:speed}. All
backends were exercised at the same \SI{15}{\minute} control protocol
on a single CPU thread, with 100 episodes of 96 transitions each
(9{,}600 total transitions per backend). Surrogate timings exclude
one-shot model loading.

The live \boptest{} RTE HTTP backend reaches \num{21.0} environment
steps per second --- the same path that downstream RL training and live
closed-loop validation use, including HTTP serialization overhead,
Flask request handling, JSON encoding, and Docker container isolation.
The control-oriented \texttt{v3} surrogate reaches \num{4626.4} steps/s
($220.2\times$ speed-up), the calibrated physical twin \texttt{v3.5}
reaches \num{2399.5} steps/s ($114.2\times$), and the canonical hybrid
backend reaches \num{1786.8} steps/s ($85.0\times$). The hybrid is
slower than \texttt{v3} alone because it evaluates both networks plus
the disagreement terms in Eq.~\eqref{eq:hybrid_loss} on every step;
$85\times$ is the practical speed-up factor relevant for the rest of
the paper, since all canonical RL training is done on the hybrid
backend.

\paragraph{Scope of the comparison.}
The $85\times$ figure compares the surrogate to the \emph{standard
\boptest{} RTE deployment} --- the HTTP--Docker stack used by the
\boptest{} community for production benchmarking and used here for all
downstream live closed-loop validation. This comparison conflates two
costs: the intrinsic Modelica/FMU simulation cost and the deployment
overhead (HTTP roundtrip, JSON serialization, Flask request handling,
container isolation). To isolate the simulation cost, we additionally
implemented a benchmark of direct FMU loading via the \texttt{fmpy}
co-simulation interface (script
\texttt{evaluation/build\_speed\_benchmark\_fmu\_direct.py}). The
\boptest{} \texttt{bestest\_air} FMU
(\texttt{boptest\_rte/testcases/bestest\_air/models/wrapped.fmu})
ships with a Linux \texttt{x86\_64} binary only
(\texttt{binaries/linux64/wrapped.so}); the archive contains no
Windows or macOS native binary, which is the reason \boptest{} is
conventionally deployed in a Linux container. Because all measurements
in this study were collected on a Windows host (matching the host of
the live HTTP benchmark for cross-comparability), the direct FMU
benchmark cannot be run in-process on the same machine without a
Linux runtime. The script is committed with explicit reproducibility
instructions for WSL2 Ubuntu and Docker, and the platform constraint
is logged at
\texttt{reports/speed\_benchmark\_fmu\_direct\_platform\_log.txt}.

We accordingly frame the reported speed-up as a comparison against the
\emph{practical \boptest{} deployment configuration} rather than
against bare FMU evaluation. A direct FMU--surrogate comparison
remains a deferred reproducibility item: when executed on a Linux
host, it will produce a third row in Table~\ref{tab:speed} that
isolates the deployment-overhead component of the current $85\times$
figure. We do not anticipate the surrogate losing its speed advantage,
because the surrogate forward pass is a fixed-cost tensor evaluation
($\sim$\SI{0.5}{\milli\second} on CPU) whereas the FMU integrates a
constrained nonlinear Modelica system per step; however, the precise
ratio against bare FMU is not currently reported and is acknowledged
as a limitation in Section~\ref{sec:discussion}.

\begin{table}[t]
  \centering
  \caption{Throughput comparison on CPU at the same \SI{15}{\minute} control protocol (100 episodes $\times$ 96 steps = 9{,}600 transitions per backend). The live \boptest{} RTE row uses the same HTTP API path that downstream RL training and live closed-loop validation actually use, so the speed-up factor in the last column reflects the practical, not the idealized, ratio. Surrogate timings exclude one-shot model loading. Source: \texttt{reports/speed\_benchmark\_table.csv}.}
  \label{tab:speed}
  \begin{tabular}{lrrrr}
    \toprule
    Backend & Steps/s & Median step (ms) & P95 step (ms) & Speed-up \\
    \midrule
    \boptest{} RTE (HTTP) & 21.0 & 14.915 & 19.251 & 1.0$\times$ \\
    \texttt{v3} surrogate & 4,626.3 & 0.165 & 0.273 & 220.2$\times$ \\
    \texttt{v3.5} calibrated & 2,399.5 & 0.359 & 0.495 & 114.2$\times$ \\
    \texttt{hybrid\_l010} & 1,786.8 & 0.507 & 0.648 & 85.0$\times$ \\
    \bottomrule
  \end{tabular}
\end{table}

